\begin{question}{Efficiently Simulating Bounded Stacks with RNNs}

\begin{subquestion}

    Result:
    \begin{table}[H]
        \begin{tabular}{cccc}
            \toprule
            & $\stackseq_1 = \openBr{2}$ &
            $\stackseq_2 = \openBr{1} \openBr{1} \openBr{1}$ &
            $\stackseq_3 = \openBr{2} \openBr{1} \openBr{2}$ \\
            \midrule
            $\closeBr{2}$ & $\emptyset$ & "reject" & $\openBr{2} \openBr{1}$ \\
            % \midrule
            $\closeBr{1}$ & "reject" & $\openBr{1} \openBr{1}$ & "reject" \\
            % \midrule
            $\openBr{2}$ & $\openBr{2} \openBr{2}$ & "reject" & "reject" \\
            \bottomrule
        \end{tabular}
    \end{table}

\end{subquestion}

\begin{subquestion}

\begin{description}

    \item[$\Sigma$:]
    $\{ \openBr{1}, ..., \openBr{k} \} \cup \{ \closeBr{1}, ..., \closeBr{k} \}$

    \item[$Q$:]
    $\{ \emptyset \} \cup \bigcup_{n=1}^{m} \{ 1, ..., k \}^n$

    \item[$\delta$:]
    $\{(q_1, \openBr{n}, q_2) | q_1 \in Q$ and $q_2 \in Q$ and $q_2 = q_1 \cdot n \} \cup \{(q_1, \closeBr{n}, q_2) | q_1 \in Q$ and $q_2 \in Q$ and $q_1 = q_2 \cdot n \}$

    \item[$I$:]
    $\{ \emptyset \}$

    \item[$F$:]
    $\{ \emptyset \}$

\end{description}

Polished:

\begin{description}
    \item[$\Sigma$:] The alphabet consisting of $k$ types of brackets: 
    $\Sigma = \{ \openBr{1}, \dots, \openBr{k} \} \cup \{ \closeBr{1}, \dots, \closeBr{k} \}$.

    \item[$Q$:] The set of states representing all valid stack sequences up to depth $m$, plus a dedicated rejection state:
    $Q = \{ s \in \{ \openBr{1}, \dots, \openBr{k} \}^* \mid |s| \le m \} \cup \{ q_{rej} \}$.
    The empty sequence $\epsilon$ represents the empty stack.

    \item[$\delta$:] The transition function $\delta: Q \times \Sigma \to Q$. For any state $q \in Q \setminus \{q_{rej}\}$:
    \begin{itemize}
        \item \textbf{Push:} For an opening bracket $\openBr{n} \in \Sigma$:
        \[ \delta(q, \openBr{n}) = 
        \begin{cases} 
        q \cdot \openBr{n} & \text{if } |q| < m \\
        q_{rej} & \text{if } |q| = m 
        \end{cases} \]
        
        \item \textbf{Pop:} For a closing bracket $\closeBr{n} \in \Sigma$:
        \[ \delta(q, \closeBr{n}) = 
        \begin{cases} 
        q' & \text{if } q = q' \cdot \openBr{n} \\
        q_{rej} & \text{otherwise}
        \end{cases} \]
    \end{itemize}
    For the rejection state, the transition is a "sink": $\delta(q_{rej}, \sigma) = q_{rej}$ for all $\sigma \in \Sigma$.

    \item[$I$:] The initial state (starting with an empty stack): $q_0 = \epsilon$.

    \item[$F$:] The set of final (accepting) states: $F = \{ \epsilon \}$.
\end{description}

\end{subquestion}

\begin{subquestion}

    Result:
    \begin{table}[H]
        \begin{tabular}{cccc}
            \toprule
            & $\stackseq_1 = \openBr{2} \begin{pmatrix} 0 \\ 1 \\ 0 \\ 0 \\ 0 \\ 0 \end{pmatrix}$ &
            $\stackseq_2 = \openBr{1} \openBr{1} \openBr{1} \begin{pmatrix} 1 \\ 0 \\ 1 \\ 0 \\ 1 \\ 0 \end{pmatrix}$ &
            $\stackseq_3 = \openBr{2} \openBr{1} \openBr{2} \begin{pmatrix} 0 \\ 1 \\ 1 \\ 0 \\ 0 \\ 1 \end{pmatrix}$ \\
            \midrule
            $\closeBr{2}$ & $\emptyset \begin{pmatrix} 0 \\ 0 \\ 0 \\ 0 \\ 0 \\ 0 \end{pmatrix}$ & "reject" & $\openBr{2} \openBr{1} \begin{pmatrix} 1 \\ 0 \\ 0 \\ 1 \\ 0 \\ 0 \end{pmatrix}$ \\
            % \midrule
            $\closeBr{1}$ & "reject" & $\openBr{1} \openBr{1} \begin{pmatrix} 1 \\ 0 \\ 1 \\ 0 \\ 0 \\ 0 \end{pmatrix}$ & "reject" \\
            % \midrule
            $\openBr{2}$ & $\openBr{2} \openBr{2} \begin{pmatrix} 0 \\ 1 \\ 0 \\ 1 \\ 0 \\ 0 \end{pmatrix}$ & "reject" & "reject" \\
            \bottomrule
        \end{tabular}
    \end{table}

\end{subquestion}

\begin{subquestion}

\[
\Sigma_{\text{open}} = \{ \langle_1, \dots, \langle_k \}, \quad \Sigma_{\text{close}} = \{ \rangle_1, \dots, \rangle_k \}
\]

\[
U(y_t) = \begin{cases} 
    U_{\text{push}} & \text{if } y_t \in \Sigma_{\text{open}} \\
    U_{\text{pop}} & \text{if } y_t \in \Sigma_{\text{close}} \\
    \mathbf{I}_{km} & \text{if } y_t = \text{EOS}
\end{cases}
\]

\[
U_{\text{push}} = \begin{pmatrix}
    \mathbf{0}_{k \times k} & \mathbf{0}_{k \times (m-1)k} \\
    \mathbf{I}_{(m-1)k} & \mathbf{0}_{(m-1)k \times k}
\end{pmatrix}, \quad
U_{\text{pop}} = \begin{pmatrix}
    \mathbf{0}_{(m-1)k \times k} & \mathbf{I}_{(m-1)k} \\
    \mathbf{0}_{k \times k} & \mathbf{0}_{k \times (m-1)k}
\end{pmatrix}
\]

\[
V = \begin{pmatrix}
    \mathbf{I}_{k} & \mathbf{0}_{k \times (k+1)} \\
    \mathbf{0}_{(m-1)k \times k} & \mathbf{0}_{(m-1)k \times (k+1)}
\end{pmatrix}
\]

\end{subquestion}

\begin{subquestion}

To simulate the bounded stack with a fixed recurrence matrix, we define a dual-track hidden state $h_t \in \mathbb{R}^{2km}$ partitioned into a top "push" track and a bottom "pop" track. At any time $t$, one track is active (encoding the stack state) while the other remains dormant (all zeros). 

Following the requirement to perform both operations simultaneously, the fixed recurrence matrix $U \in \mathbb{R}^{2km \times 2km}$ is defined to draw from either track of $h_{t-1}$ to generate candidates for both tracks in the next state:
\[
U = \begin{pmatrix}
    U_{\text{push}} & U_{\text{push}} \\
    U_{\text{pop}} & U_{\text{pop}}
\end{pmatrix}
\]
where $U_{\text{push}}$ and $U_{\text{pop}}$ are the $km \times km$ transition matrices from Question 3d. To implement selection logic via the Heaviside activation we define a global bias vector $b \in \mathbb{R}^{2km}$ that "locks" the state by default:
\[
b = -\mathbf{1}_{2km} = \begin{pmatrix} 
    -1 \\ 
    \vdots \\ 
    -1 
\end{pmatrix}
\]
The input matrix $V \in \mathbb{R}^{2km \times (2k+1)}$ injects bracket identities and provides a $+1$ gating signal to counteract the bias for the correct track. Following the required symbol order:
\[
V = \begin{pmatrix}
    \mathbf{I}_{k} + \mathbf{1}_{k \times k} & \mathbf{0}_{k \times k} & \mathbf{1}_{k \times 1} \\
    \mathbf{1}_{(m-1)k \times k} & \mathbf{0}_{(m-1)k \times k} & \mathbf{1}_{(m-1)k \times 1} \\
    \mathbf{0}_{k \times k} & \mathbf{1}_{k \times k} & \mathbf{1}_{k \times 1} \\
    \mathbf{0}_{(m-1)k \times k} & \mathbf{1}_{(m-1)k \times k} & \mathbf{1}_{(m-1)k \times 1}
\end{pmatrix}
\]
The hidden state update $h_t = H(Uh_{t-1} + Vo_{y_t} + b)$ ensures that if $y_t$ is an opening bracket, the top track receives a $+1$ boost from $V$ to unlock the "push" operation ($1 + 1 - 1 = 1$), while the dormant bottom track is suppressed to $0$. A closing bracket similarly unlocks the bottom track, while EOS provides a boost to both tracks to preserve the final configuration.

-----

To prove the RNN simulates the FSA $\mathcal{A}$, we define an embedding function $\phi: Q \to \mathbb{R}^{2km}$ such that for a stack sequence $\gamma$, $\phi([\gamma])$ corresponds to the encoding $h(\gamma)$ residing in exactly one half of the hidden state while the other half is zeroed. We proceed by induction on the string length $t$. For the base case, the empty stack corresponds to $h_0 = \mathbf{0}_{2km}$. 

For the inductive step, assume $h_{t-1}$ correctly encodes the stack state $q_{t-1}$. Due to the block structure of $U$, the product $U h_{t-1}$ generates a "dual candidate" vector:
\[
U h_{t-1} = \begin{pmatrix} U_{\text{push}} \\ U_{\text{pop}} \end{pmatrix} (\text{active half of } h_{t-1}) = \begin{pmatrix} h_{\text{push candidate}} \\ h_{\text{pop candidate}} \end{pmatrix}
\]
The behavior is then determined by the input $y_t$ \cite{286}:
\begin{enumerate}
    \item \textbf{Push Case ($y_t = \langle_i$):} The input matrix $V$ provides a $+1$ boost and the symbol identity $e_i$ to the top track. In the bottom track, the bias $b=-\mathbf{1}$ remains uncompensated. Applying $H(x > 0)$ yields:
    \[
    h_t = \begin{pmatrix} H(h_{\text{push candidate}} + e_i + 1 - 1) \\ H(h_{\text{pop candidate}} + 0 - 1) \end{pmatrix} = \begin{pmatrix} h(\gamma_t) \\ \mathbf{0}_{km} \end{pmatrix}
    \]
    \item \textbf{Pop Case ($y_t = \rangle_i$):} $V$ provides the $+1$ boost only to the bottom track. The top track is suppressed by the bias, while the bottom track survives:
    \[
    h_t = \begin{pmatrix} H(h_{\text{push candidate}} + 0 - 1) \\ H(h_{\text{pop candidate}} + 1 - 1) \end{pmatrix} = \begin{pmatrix} \mathbf{0}_{km} \\ h(\gamma_t) \end{pmatrix}
    \]
\end{enumerate}
In both cases, $h_t$ correctly represents the new FSA state $q_t = \delta(q_{t-1}, y_t)$ in either the top or bottom track \cite{278}. This induction shows that the recurrence relation $h_t = H(Uh_{t-1} + Vo_{y_t} + b)$ is equivalent to the FSA transition function $\delta$, successfully simulating a bounded stack with a fixed recurrence matrix \cite{286}.

\end{subquestion}

\begin{subquestion}

To decode the $2km$-dimensional hidden state from 3e using the encoding in Eq. (27), we define $E$ as a redundant block matrix:
\[ E = \begin{pmatrix} E_{\text{track}} & E_{\text{track}} \end{pmatrix} \in \mathbb{R}^{(2k+1) \times 2km} \]
where $E_{\text{track}} \in \mathbb{R}^{(2k+1) \times km}$ and $b' \in \mathbb{R}^{2k+1}$ are defined using parameter $\alpha > 0$:
\[
b' = \begin{pmatrix} \alpha \cdot \mathbf{1}_k \\ \mathbf{0}_k \\ \alpha \end{pmatrix}, \quad 
E_{\text{track}} = \begin{pmatrix} 
\mathbf{0}_{k \times (m-1)k} & -\alpha \cdot \mathbf{1}_{k \times k} \\
\alpha \cdot \mathbf{I}_k & \mathbf{0}_{k \times (m-1)k} \\
-\alpha \cdot \mathbf{1}_{1 \times mk}
\end{pmatrix}
\]

\paragraph{Logit Logic for Recognition}
The construction ensures logits $z_j = \alpha$ for valid symbols and $z_j \le 0$ for invalid ones:
\begin{itemize}
    \item \textbf{Opening Brackets:} Use \textit{Veto Logic}. $b'$ grants default permission ($\alpha$), but $E_{\text{track}}$ subtracts $\alpha$ if the $m$-th stack slot is occupied, rejecting the symbol when the depth limit is reached.
    \item \textbf{Closing Brackets:} Use \textit{Permission Logic}. $b'$ is $0$ (default rejection). $E_{\text{track}}$ adds $\alpha$ only if the corresponding opening bracket $i$ matches the stack top (Slot 1).
    \item \textbf{EOS:} Uses \textit{Veto Logic}. $b'$ is $\alpha$, but $E_{\text{track}}$ subtracts $\alpha$ if any bracket exists in the stack, ensuring termination only for empty stacks.
\end{itemize}

\paragraph{Bounding $\alpha$}
To satisfy $\eta$-truncated support ($p \ge \eta$), we consider the normalization constant $Z = \sum \exp(z_j)$. The worst-case (lowest probability) occurs at empty or partial stack depths where $k+1$ symbols are valid:
\[ Z_{\text{max}} = (k+1)e^{\alpha} + k \cdot e^0 \]
We require the probability of a valid symbol to exceed $\eta$:
\[ \frac{e^{\alpha}}{(k+1)e^{\alpha} + k} \ge \eta \implies e^{\alpha}(1 - \eta(k+1)) \ge \eta k \]
Given $\eta < \frac{1}{k+1}$, we set the lower bound for $\alpha$ as:
\[ \alpha \ge \ln \left( \frac{\eta k}{1 - \eta(k+1)} \right) \]
This ensures $\mathcal{L}_{\eta}(p_{\text{SM}}) = D(k,m)$ by squeezing all invalid continuations below the threshold.

\end{subquestion}

\begin{subquestion}

\end{subquestion}

\end{question}